% Related Work — claude-tunnel bachelor thesis.
\chapter{Verwandte Arbeiten}
\label{ch:relatedwork}

Der Entwurf beruehrt drei Forschungsstraenge: verschluesselte Tunnel und
Anonymitaetsnetze, providerblinde Vermittlung sowie Zensurumgehung. Dieses
Kapitel ordnet die vorliegende Arbeit in jeden dieser Straenge ein und arbeitet
heraus, worin die Kombination neu ist.

\section{Verschluesselte Tunnel und VPNs}
\label{sec:rw-vpn}

Klassische \gls{ac:vpn}-Loesungen bauen einen verschluesselten Tunnel zwischen
Client und einem Konzentrator auf. WireGuard~\cite{donenfeld2017wireguard} ist
hierfuer ein modernes, schlankes Beispiel: es nutzt -- wie die vorliegende
Arbeit -- Bausteine aus dem Noise-Umfeld (Curve25519, ChaCha20-Poly1305,
BLAKE2s) und erreicht einen kompakten, gut analysierbaren Handshake. Der
entscheidende Unterschied liegt im Vertrauensmodell: der VPN-Server ist der
Endpunkt des Tunnels und sieht damit die Ziele des Clients im Klartext. Die
Vertraulichkeit gegenueber dem Betreiber ist keine strukturelle Eigenschaft,
sondern beruht auf dessen Wohlverhalten. Die vorliegende Arbeit uebernimmt die
kryptographische Substanz, verschiebt den Terminierungspunkt der
Nutzlast-Verschluesselung aber vom Vermittler zum Origin.

\section{Anonymitaetsnetze}
\label{sec:rw-tor}

Tor~\cite{dingledine2004tor} entkoppelt Absender und Ziel ueber eine Kette von
mindestens drei Relays, von denen keines beide Enden zugleich kennt. Das
Schutzziel ist primaer die \emph{Unverkettbarkeit} von Nutzer und Ziel gegenueber
einzelnen Relays, erkauft durch mehrere Zwischensprunge und entsprechende
Latenz. Die vorliegende Arbeit verfolgt ein engeres, aber schaerferes Ziel: die
strukturelle Blindheit eines \emph{einzelnen} Vermittlers gegenueber
\emph{Inhalt und Ziel}, ergaenzt um einen direkten Pfad, der diesen Vermittler
bei Erreichbarkeit ganz umgeht. Anonymitaet im Tor'schen Sinne ist ausdruecklich
kein Ziel; wohl aber die Nicht-Einsehbarkeit durch den Betreiber.

\section{Providerblinde Vermittlung}
\label{sec:rw-ohttp}

Am naechsten steht die Arbeit dem Prinzip von Oblivious
HTTP~\cite{rfc9458}: dort trennt ein Relay das \emph{Wer} vom \emph{Was}, indem
es Anfragen an ein Gateway weiterreicht, ohne deren Inhalt zu sehen, waehrend das
Gateway den Absender nicht kennt. Oblivious HTTP ist jedoch auf einzelne,
kapselbare Anfrage-Antwort-Paare zugeschnitten und setzt ein kooperierendes
Gateway voraus. Die vorliegende Arbeit uebertraegt den Grundgedanken -- der
Vermittler sieht nur Chiffretext -- auf einen \emph{allgemeinen, bidirektionalen
Tunnel} beliebiger \gls{ac:tcp}- und \gls{ac:udp}-Protokolle, mit
Origin-Pinning statt Gateway und mit einem Proof-of-Work-Rendezvous als
Missbrauchsschranke.

\section{Proxying ueber HTTP und QUIC}
\label{sec:rw-masque}

MASQUE (\enquote{Proxying UDP in HTTP})~\cite{rfc9298} tunnelt \gls{ac:udp}
ueber HTTP/3 und damit ueber \gls{ac:quic}. Technisch aehnelt das dem hier
genutzten Transport, das Vertrauensmodell unterscheidet sich jedoch erneut: der
MASQUE-Proxy kennt das Ziel (Adresse und Port), weil er selbst die
UDP-Assoziation aufbaut. Er ist ein transparenter Weiterleiter, kein blinder
Relay. Der Fallback der vorliegenden Arbeit (\gls{ac:tls}-ueber-\gls{ac:tcp})
verfolgt ein verwandtes Ziel wie MASQUE -- Betrieb, wenn \gls{ac:udp} blockiert
ist --, transportiert dabei aber weiterhin nur den bereits Ende-zu-Ende
verschluesselten Tunnel.

\section{Zensurumgehung}
\label{sec:rw-censorship}

Zur Umgehung aktiver Netzsperren existiert ein breites Feld. Domain
Fronting~\cite{fifield2015domain} verbirgt das wahre Ziel hinter dem Namen eines
grossen, schwer blockierbaren Dienstes; Decoy Routing bzw.\ Refraction
Networking~\cite{wustrow2011telex} verlagert die Umlenkung in die
Netzinfrastruktur selbst. Diese Ansaetze zielen primaer auf die
\emph{Unauffaelligkeit} des Datenstroms (Verschleierung gegen einen
beobachtenden Zensor). Die vorliegende Arbeit setzt komplementaer an: nicht
Verschleierung, sondern strukturelle Blindheit des Vermittlers plus ein
Transport-Fallback, der die haeufigste grobe Sperre -- das Blockieren von
\gls{ac:udp} -- ueberbrueckt. Traffic-Analyse-Resistenz gegen einen aktiven
Beobachter ist kein Ziel und wird in Kapitel~\ref{ch:evaluation} als offene
Grenze benannt.

\section{Einordnung}
\label{sec:rw-einordnung}

Die genannten Systeme decken jeweils einen Teil der Anforderungen ab:
WireGuard die kryptographische Substanz, Oblivious HTTP die Blindheit, MASQUE den
QUIC-Transport, Tor die Entkopplung, die Zensurumgehungs-Systeme die
Sperrresistenz. Tabelle~\ref{tab:rw-vergleich} stellt die Systeme entlang der
zentralen Eigenschaften gegenueber und macht sichtbar, dass keines sie
vollstaendig vereint.

\begin{table}[t]
  \centering
  \caption{Verwandte Systeme entlang der zentralen Eigenschaften.
  $\bullet$ = erfuellt, $\circ$ = teilweise, -- = nicht erfuellt.}
  \label{tab:rw-vergleich}
  \begin{tabular}{lccccc}
    \toprule
    System & E2E, Vermittler-blind & Allg.\ TCP/UDP & QUIC & P2P-Pfad & Missbrauchsschranke \\
    \midrule
    WireGuard~\cite{donenfeld2017wireguard} & $\circ$ & $\bullet$ & -- & $\bullet$ & -- \\
    Tor~\cite{dingledine2004tor} & $\circ$ & $\circ$ & -- & -- & -- \\
    Oblivious HTTP~\cite{rfc9458} & $\bullet$ & -- & -- & -- & -- \\
    MASQUE~\cite{rfc9298} & -- & $\circ$ & $\bullet$ & -- & -- \\
    Diese Arbeit & $\bullet$ & $\bullet$ & $\bullet$ & $\bullet$ & $\bullet$ \\
    \bottomrule
  \end{tabular}
\end{table}

Die Zeilen bestaetigen die Prosa: WireGuard verschluesselt zwar Ende-zu-Ende,
doch jeder Peer kennt seinen Gegenpart ($\circ$); Tor entkoppelt zwar Absender
und Ziel, ist aber auf \gls{ac:tcp} zugeschnitten ($\circ$) und ohne
Direktpfad; Oblivious HTTP ist blind, aber auf Anfrage-Antwort-Paare begrenzt;
MASQUE bringt den QUIC-Transport, sein Proxy ist jedoch sehend. Der Beitrag der
vorliegenden Arbeit ist die \emph{Kombination}
dieser Eigenschaften in einem einzigen, lauffaehigen System: ein providerblinder,
allgemein einsetzbarer Tunnel mit Ende-zu-Ende-Pinning, Proof-of-Work-Rendezvous,
direktem Peer-to-Peer-Pfad und Transport-Fallback -- inklusive empirischer
Vermessung des dadurch entstehenden Latenz-Overheads.
